\documentclass[../../../include/open-logic-section]{subfiles}

\begin{document}

\olfileid{his}{set}{pathology}
\olsection{विलक्षण रचना}

मात्र \emph{सीमेच्या} व्याख्येमुळे काही खरोखर “विलक्षण” रचना शक्य असल्याचे दिसून आले.

१८३०च्या सुमारास बोल्झानो यांनी असे एक फलन शोधले जे \emph{सर्वत्र सतत} होते, पण \emph{कोठेही अवकलनीय नव्हते}. (दुर्दैवाने बोल्झानो यांनी हे कधी प्रकाशित केले नाही. वायेरश्ट्रास यांनी स्वतंत्रपणे हीच कल्पना शोधल्यामुळे गणितज्ञांना १८७२ मध्ये तिची प्रथम ओळख झाली.)\footnote{हा इतिहास \href{http://en.wikipedia.org/wiki/Weierstrass_function}{वायेरश्ट्रास फलनावरील विकिपीडिया लेखाच्या} अत्यंत तपशीलवार तळटीपांत नोंदवला आहे.} हे किमान आश्चर्यकारक तरी होतेच. $|x|$ सारखी सर्वत्र सतत, पण एखाद्या विशिष्ट बिंदूवर अवकलनीय नसलेली फलने सहज सापडतात. परंतु सर्वत्र सतत असूनही \emph{कोठेही} अवकलनीय नसलेले फलन अगदी वेगळेच आहे. असे फलन कसे काढाल याचा क्षणभर विचार करा. ते सतत असण्यासाठी कागदावरून लेखणी कधीही न उचलता त्याचा आलेख काढता आला पाहिजे; पण ते कोठेही अवकलनीय नसण्यासाठी लेखणीची दिशा सतत अचानक बदलावी लागेल.

यानंतर आणखी “विलक्षण” रचना समोर आल्या. ५ जानेवारी १८७४ रोजी कँटर यांनी डेडेकिंड यांना पत्र लिहून पुढील प्रश्न विचारला:
\begin{quote}
एखाद्या पृष्ठभागाची (उदाहरणार्थ सीमेसह चौरसाची) एखाद्या रेषेशी (उदाहरणार्थ टोकांसह सरळ रेषेशी) अशी एकास-एक संगती लावता येईल का की पृष्ठभागावरील प्रत्येक बिंदूला रेषेवरील एक बिंदू मिळेल आणि उलट रेषेवरील प्रत्येक बिंदूला पृष्ठभागावरील एक बिंदू मिळेल?

या प्रश्नाचे उत्तर शोधणे अजूनही मला फार कठीण वाटते—तरी येथेही \emph{नाही} असे म्हणण्याची इतकी तीव्र ओढ वाटते की त्यासाठी सिद्धतेची गरज जवळजवळ नसावी असे वाटू लागते.
[\citealt{Gouvea2011} मधून उद्धृत]
\end{quote}
पण १८७७ मध्ये कँटर यांनी स्वतःचा अंदाज चुकीचा असल्याचे सिद्ध केले. प्रत्यक्षात रेषा आणि चौरस यांवरील बिंदूंची संख्या अगदी सारखी असते. २९ जून १८७७ रोजी त्यांनी डेडेकिंड यांना “\emph{je le vois, mais je ne le crois pas}” असे लिहिले; म्हणजे, “ते मला दिसते, पण त्यावर माझा विश्वास बसत नाही.” गणिताच्या “प्रचलित इतिहासकथनात” याचा अर्थ अनेकदा असा घेतला जातो की त्या काळच्या गणितज्ञांना हे नवे निष्कर्ष \emph{अक्षरशः अविश्वसनीय} वाटत होते. (हा पत्रव्यवहार \citet{Gouvea2011} येथे दिला आहे आणि \olref[his][set][mythology]{sec} मध्ये आपण त्याकडे पुन्हा वळू. कँटर यांच्या सिद्धतेचा आराखडा \olref[his][set][cantorplane]{sec} मध्ये दिला आहे.)

कँटर यांच्या निष्कर्षाने प्रेरित होऊन, एखाद्या रेषेचे \emph{अखंडपणे} प्रतलावर प्रतिचित्रण करून ते प्रतल भरून काढता येईल का, असा विचार पियानो यांनी सुरू केला. असे प्रतिचित्रण म्हणजे \emph{अवकाश भरणारा वक्र} असेल. \citeyear{Peano1890} मध्ये पियानो यांनी असाच वक्र रचला. हे अंतर्ज्ञानाला खरोखर विरुद्ध वाटते: युक्लिडने रेषेची व्याख्या “रुंदीविरहित लांबी” अशी केली होती (पुस्तक I, व्याख्या २); पण रेषा योग्य प्रकारे वळवली, तर तिची लांबी रुंदी व्यापू शकते, हे पियानो यांनी दाखवले. \citeyear{Hilbert1891} मध्ये हिल्बर्ट यांनी याहून किंचित अधिक सहज समजणारा अवकाश भरणारा वक्र आणि त्याची काही चित्रे दिली. हा वक्र टप्प्याटप्प्याने रचला जातो; रचनेचे पहिले सहा टप्पे येथे दिले आहेत:
\begin{center}
\begin{tikzpicture}
\begin{scope}[xshift=20pt, yshift=20pt]
	\draw [oldiagcolorC, l-system={Hilbert curve, step=40pt, angle=90, axiom=L, order=1}] lindenmayer system;
\end{scope}
\begin{scope}[xshift=110pt, yshift=10pt]
	\draw [oldiagcolorC, l-system={Hilbert curve, step=20pt, angle=90, axiom=L, order=2}] lindenmayer system;
\end{scope}
\begin{scope}[xshift=205pt, yshift=5pt]
	\draw [oldiagcolorC, l-system={Hilbert curve, step=10pt, angle=90, axiom=L, order=3}] lindenmayer system;
\end{scope}
\begin{scope}[xshift=2.5pt, yshift=-97.5pt]
	\draw [oldiagcolorC, l-system={Hilbert curve, step=5pt, angle=90, axiom=L, order=4}] lindenmayer system;
\end{scope}
\begin{scope}[xshift=101.25pt, yshift=-98.75pt]
	\draw [oldiagcolorC, l-system={Hilbert curve, step=2.5pt, angle=90, axiom=L, order=5}] lindenmayer system;
\end{scope}
\begin{scope}[xshift=200.625pt, yshift=-99.375pt]
	\draw [oldiagcolorC, l-system={Hilbert curve, step=1.25pt, angle=90, axiom=L, order=6}] lindenmayer system;
\end{scope}
\draw[gray] (0pt,0pt) rectangle (80pt, 80pt);
\draw[gray] (100pt,0pt) rectangle (180pt, 80pt);
\draw[gray] (200pt,0pt) rectangle (280pt, 80pt);
\draw[gray] (0pt,-100pt) rectangle (80pt, -20pt);
\draw[gray] (100pt,-100pt) rectangle (180pt, -20pt);
\draw[gray] (200pt,-100pt) rectangle (280pt, -20pt);
\end{tikzpicture}
\end{center}
सीमेत—जी संकल्पना आता काटेकोरपणे व्याख्यित झाली होती—संपूर्ण चौरस एकसंध !!{colorC} रंगाने भरतो. तसेच, हिल्बर्ट यांचा वक्र सर्वत्र सतत आहे, पण कोठेही अवकलनीय नाही; अंतर्ज्ञानाने सांगायचे तर, अनंत टप्प्यांच्या सीमेत फलन प्रत्येक क्षणी अचानक दिशा बदलते. (हिल्बर्ट यांच्या रचनेचा अधिक तपशीलवार आराखडा \olref[his][set][hilbertcurve]{sec} मध्ये पाहू.)

या “विलक्षण” भूमितीय रचनांना, योग्य की अयोग्य, भूमितीय अंतर्ज्ञानावर विसंबून राहण्याबाबत शंका घेण्याचे कारण मानले गेले. गणित करण्याच्या एका नव्या पद्धतीच्या \emph{प्रचाराचे} त्या जणू साधन बनल्या; या पद्धतीचा विकास पुढे संच उपपत्तीत झाला. या विषयाभोवती नंतर उभ्या राहिलेल्या दंतकथांमध्ये, हे निष्कर्ष पूर्णपणे काटेकोर आणि तितकेच धक्कादायक होते, असे वारंवार सांगितले गेले. त्यामुळे त्यांनी दोन कामे केली: भूमितीय अंतर्ज्ञानावर अवलंबून राहण्याविरुद्ध सावध केले आणि नव्या विचारपद्धतींची फलद्रूपता दाखवली.

\end{document}
